\stitle{Related work}. We categorized the related work as follows.

\etitle{Graph pattern matching algorithm}.
Algorithms for subgraph isomorphism have been
extensively studied.
Ullmann proposed the first backtracking-based algorithm for subgraph isomorphism~\cite{Iso-match-6}. 
Subsequent research in algorithms for exact graph pattern 
matching focuses on pruning methods~\cite{carletti2017introducing,carletti2019vf3}, 
the application of index~\cite{CECI,sun2020subgraph}, 
index methods based on data mining~\cite{yan2004graph,zhao2007graph}, and so on. 
\looseness=-1

Isomorphic mappings suffers from two major drawbacks:
(1) the decision version of subgraph isomorphism is 
NP-complete~\cite{cook2023complexity},
(2) The requirement of  
\revise{injective} map 
could be too strict in many cases,
especially when the data source is noisy~\cite{yuan2011efficient}. 
 %
To address these challenges, 
(1) a series of studies focus on designing similarity functions and 
deriving matching results through  
top-k candidates or the establishment of a similarity threshold
\cite{fan2010graph,cheng2013top,dutta2017neighbor,wu2013ontology,khan2013nema,zhao2013partition}. 
(2) The other studies relax the  
\revise{injective} map to less restrictive  
relations.
Bounded simulation 
is developed to bound the lengths of 
matching paths~\cite{fan2010grapha}.
\revise{Dual simulation and } 
Strong simulation \revise{are proposed}
 in \cite{ma2011capturing}  to 
\revise{bound directions of edges 
	and 
the distance of vertices in a match, 
respectively}. 
Subsequently, various simulation-based
matching algorithms have been developed~\cite{fan2013incremental,fan2014distributed, fard2014distributed,gao2016prs,du2018personalized}.
\looseness = -1


\etitle{ML Models for graph pattern matching}.
Multiple ML models are developed for graph pattern matching~\cite{NeuroMatch,Sub-ML-1,Sub-ML-2,Sub-ML-3,Sub-ML-4}.
Most of these models are based on Graph Neural Networks (GNNs),
to compute embeddings of vertices in a graph~\cite{scarselli2008graph},
and make a prediction using the order-embedding space~\cite{NeuroMatch}
or MLP~\cite{Sub-ML-1}.
Most work focus on 
enhancing GNNs architecture~\cite{wang2022twin,wang2022equivariant,zhang2021eigen,bouritsas2022improving,feng2022powerful},
to learn more topological information. 
Zhang et al. \cite{zhang2023expressive} 
encoded distance information into vertex embeddings. 
Geerts et al. \cite{geerts2022expressiveness} 
proposed a model to aggregate information 
from $k$-order neighbors using computational operators. 
%
You et al. \cite{you2021identity} injected a unique 
color identity to nodes, 
and performed heterogeneous message passing 
to obtain the node embedding. 
Xu et al. \cite{xu2018powerful} 
optimized the aggregation function 
by using multisets, 
to generate distinct embeddings for different nodes. 
Maron et al. \cite{maron2018invariant} 
proposed a model consisting of 
both equivariant linear layers and nonlinear activation layers. 

There has also been work on
studying properties of GNNs.
(1) The expressive power of these GNN
models can be bounded by 
the Weisfeiler-Lehman (WL) hierarchy~\cite{xu2018powerful,morris2019weisfeiler,barcelo2022weisfeiler,zhang2023expressive,GNN-WL-C},
descriptive complexity~\cite{Express-LICS},
fragments of FO logic~\cite{Expressive-FOC2}
and global feature map transformer~\cite{Express-GFMT}. 
See~\cite{zhang2023expressive} for a recent survey.
(2) The generalization error
bound for GNNs can be measured by 
Rademacher complexity~\cite{Gen-bound}, VC dimension~\cite{VC-GNN}
and generalization gap~\cite{Generalization-gap,Generalization-gap}.
